\documentclass[11pt]{article}
\usepackage[T1]{fontenc}
\usepackage{lmodern,amsmath,amssymb,amsthm,mathtools}
\usepackage[margin=1in]{geometry}
\usepackage{booktabs,microtype,enumitem,xcolor,fancyhdr}
\usepackage[numbers,sort&compress]{natbib}
\usepackage{hyperref}
\hypersetup{colorlinks=true,linkcolor=blue!45!black,citecolor=blue!45!black,
  urlcolor=blue!45!black,pdftitle={An upper bound on the Berry-Esseen constant for independent summands},
  pdfauthor={Logan Bell}}
\setlength{\emergencystretch}{2em}
\setlength{\headheight}{14pt}
\pagestyle{fancy}
\fancyhf{}
\fancyhead[L]{\small The independent Berry--Esseen constant}
\fancyhead[R]{\small Research draft}
\fancyfoot[C]{\thepage}
\numberwithin{equation}{section}
\newtheorem{candidate}{Proposed theorem}[section]
\newtheorem{lemma}[candidate]{Lemma}
\newtheorem{proposition}[candidate]{Proposition}
\theoremstyle{remark}
\newtheorem{remark}[candidate]{Remark}
\DeclareMathOperator{\Var}{Var}
\DeclareMathOperator{\sinc}{sinc}
\DeclareMathOperator{\Unif}{Unif}
\newcommand{\E}{\mathbb E}
\newcommand{\Prob}{\mathbb P}
\newcommand{\R}{\mathbb R}
\newcommand{\CE}{C_{\mathrm E}}
\newcommand{\Cind}{C_{\mathrm{ind}}}
\newcommand{\ii}{\mathrm i}
% Generated from result.json by verify/paper_values.py.
\newcommand{\UpperBound}{0.44995}
\newcommand{\ExactUpper}{\frac{8105578597689285}{18014398509481984}}
\newcommand{\ReplayedCount}{141,533}
\newcommand{\SharedBoundCount}{92,041}
\newcommand{\SmallUpper}{0.447913037296}
\newcommand{\VarianceUpper}{0.540936541549}
\newcommand{\LargeUpper}{0.447055000001}
\newcommand{\BandCount}{1,670}
\newcommand{\NodeCount}{478,804}
\newcommand{\AcceptedCount}{233,574}
\newcommand{\InfeasibleCount}{6,663}
\newcommand{\RecordedUpper}{0.44994999935317664}


\title{An upper bound on the Berry--Esseen constant\\ for independent summands}
\author{Logan Bell}
\date{7 September 2026}

\begin{document}
\maketitle
\thispagestyle{fancy}
\begin{abstract}
Let $\Delta$ be the Kolmogorov distance between a standardized sum of
independent centered real random variables and the standard normal law,
and let $L$ be the normalized sum of third absolute moments. We propose
the uniform bound $\Delta<\CandidateUpper\,L$. The argument treats a
summand of maximal variance separately and bounds the real and imaginary
parts of the zero-bias error in terms of third-moment excess over
symmetric two-point laws. A quadratic term in the logarithm strengthens
the characteristic-function product bound. Prawitz's smoothing
inequality then reduces the intermediate range of $L$ to a finite
interval calculation; analytic estimates cover small and large $L$.
The supplement contains the arithmetic implementation and complete
partition records. Independent mathematical review remains outstanding.

\end{abstract}

\section{The problem and the candidate bound}\label{sec:statement}



Let $X_1,\ldots,X_n$ be a finite independent family of centered real random
variables. Assume
\[
  0<V=\sum_{j=1}^n\E X_j^2<\infty,
  \qquad \sum_{j=1}^n\E|X_j|^3<\infty.
\]
Write $\Phi$ and $\phi$ for the standard normal distribution function and
density, and set
\begin{equation}\label{eq:problem}
  S=\frac{\sum_jX_j}{\sqrt V},\qquad
  L=\frac{\sum_j\E|X_j|^3}{V^{3/2}},\qquad
  \Delta=\sup_{x\in\R}|\Prob(S\le x)-\Phi(x)|.
\end{equation}
Thus $\Cind=\sup\Delta/L$, over all such arrays, is the unrestricted
independent Berry--Esseen constant. Esseen's lower bound is
\begin{equation}\label{eq:esseen}
  \CE=\frac{\sqrt{10}+3}{6\sqrt{2\pi}}
      =0.40973218370239634299\ldots\le\Cind;
\end{equation}
see \citet{Esseen1956} and the discussion in \citet{Shevtsova2012}.
The conjectured equality $\Cind=\CE$ remains open.

\begin{candidate}\label{thm:candidate}
For the arrays in \eqref{eq:problem},
\begin{equation}\label{eq:candidate}
  \Delta<\CandidateUpper\,L<0.475\,L,
  \qquad \CE\le\Cind<\CandidateUpper.
\end{equation}
\end{candidate}

The arithmetic assumptions and completed checks for this proposed bound
are stated in Section~\ref{sec:certificate}. For comparison,
\citet{Shevtsova2013} obtained the upper bounds $0.5583$ for independent
summands and $0.4690$ for identically distributed summands.
\citet{Schulz2016} proved the sharp constant $\CE$ for binomial laws.

The proof uses the zero-bias and Fourier framework of \citet{Tyurin2009}
and the smoothing inequality recorded in \citet{Shevtsova2012}.
Comparison with symmetric two-point sums also appears in the
smooth-function results of \citet{MattnerShevtsova2019}. We derive below
the estimates used here, including their aggregation across unequal
variances.

\subsection{Normalization and moment parameters}
Replace each $X_j$ by $X_j/\sqrt V$, reuse the notation $X_j$, and
discard summands of variance zero. Then $S=\sum_jX_j$ and the total
variance is one. Put
\[
 v_j=\E X_j^2,\qquad b_j=\E|X_j|^3,\qquad
 \tau=\sum_jv_j^{3/2},\qquad \sum_jv_j=1,\quad \sum_jb_j=L.
\]
The differences $b_j-v_j^{3/2}$ are nonnegative and sum to $L-\tau$.
They vanish for centered symmetric two-point summands. Relabel the
summands so that $v_1=\max_jv_j$, and put $d=v_1$ and $b=b_1$.
The moment parameters satisfy
\begin{equation}\label{eq:moment-domain}
 b\ge d^{3/2},\qquad
 d^{3/2}\le\tau\le\min\{L-b+d^{3/2},\sqrt d\},\qquad \tau\le L.
\end{equation}
Indeed $b_j\ge v_j^{3/2}$ and $v_j\le d$ imply
$\tau\le\sqrt d\sum_jv_j=\sqrt d$. In particular
$\tau^2\le d\le\tau^{2/3}$.
For the remaining summands, write
\begin{equation}\label{eq:remainder-state}
 w=1-d,\qquad B_R=L-b,\qquad \tau_R=\tau-d^{3/2},\qquad
 a_R=\min\{d,w,B_R^{2/3},\tau_R^{2/3}\}.
\end{equation}
For $j\ge2$, we have $v_j\le d$, $v_j\le w$,
$v_j^{3/2}\le b_j\le B_R$, and $v_j^{3/2}\le\tau_R$. Thus
$v_j\le a_R$. Every realizable tuple satisfies these constraints;
maximizing over the possibly larger set they define gives an upper bound.

Section~\ref{sec:products} bounds products of characteristic functions.
Sections~\ref{sec:directional} and~\ref{sec:quadratic} estimate the
zero-bias error of each summand and aggregate these errors.
Section~\ref{sec:distinguished} combines the estimates while treating
$X_1$ separately. Smoothing and the interval calculation occupy
Sections~\ref{sec:smoothing} and~\ref{sec:certificate}.

\section{Bounds for characteristic-function products}\label{sec:products}

For a random variable $X$, let $f_X(t)=\E e^{\ii tX}$. Since
$f_X(-t)=\overline{f_X(t)}$, it suffices to bound its modulus for
$t\ge0$. For a nonempty index set $J$, put
\[
 V_0=\sum_{j\in J}v_j,\qquad B=\sum_{j\in J}b_j,\qquad
 \tau_0=\sum_{j\in J}v_j^{3/2},
\]
and let $a$ satisfy $v_j\le a$ for every $j\in J$. We need bounds for
the product over $J$ and for the product after any one index is omitted.
Products over the empty set equal one.

\subsection{Cubic and trigonometric bounds}
Define
\begin{equation}\label{eq:kappa}
 \kappa_* = \sup_{x>0}\frac{\cos x-1+x^2/2}{x^3},
 \qquad \kappa=\frac{49581}{500000}=0.099162.
\end{equation}
Let $\theta$ be the maximizer in \eqref{eq:kappa}. For $V_0,E_0>0$,
$t\ge0$, and $u=E_0t/V_0$, put
\begin{equation}\label{eq:tyurin-exponent}
 \psi(t;V_0,E_0)=
 \begin{cases}
   V_0t^2/2-\kappa_*E_0t^3,&u<\theta,\\
   V_0^3E_0^{-2}(1-\cos u),&\theta\le u\le2\pi,\\
   0,&u>2\pi.
 \end{cases}
\end{equation}
The characteristic-function bounds in \citet[Lemma~1]{Tyurin2009} give,
with $v=\E X^2>0$ and $b=\E|X|^3\ge v^{3/2}$ in the second formula,
\begin{align}
 \Big|\prod_j f_{X_j}(t)\Big|
   &\le \exp\{-\psi(t;V_0,B+\tau_0)\},\label{eq:aggregate-cf}\\
 |f_X(t)|&\le A(t;v,b)
   :=\sqrt{1-2\psi(t;v,b+v^{3/2})}.\label{eq:single-cf}
\end{align}
For fixed parameters, $\psi$ first increases and then decreases with
$t$, and is zero after the cosine branch. On the cubic branch,
$\partial_t\psi=t(V_0-3\kappa_*E_0t)$; on the cosine branch,
$\partial_t\psi=(V_0^2/E_0)\sin u\le0$. Thus the upper bounds in
\eqref{eq:aggregate-cf}--\eqref{eq:single-cf} attain their maxima on a
frequency interval at an endpoint. On the cosine branch,
\begin{align*}
 \partial_{V_0}\psi
 &=\frac{V_0^2}{E_0^2}\{3(1-\cos u)-u\sin u\}\ge0,\\
 \partial_{E_0}\psi
 &=\frac{V_0^3}{E_0^3}\{u\sin u-2(1-\cos u)\}\le0.
\end{align*}
The same signs hold on the cubic branch. The maximizing equation for
$\theta$ makes the formulas agree at the joins. Consequently the upper
bounds decrease with the variance argument and increase with the moment
argument on the domains just specified.

For $r(x)=(\cos x-1+x^2/2)/x^3$, differentiation gives
\[
 r'(x)=-\frac{h(x)}{2x^4},\quad
 h(x)=x^2+2x\sin x+6(\cos x-1),\quad
 h'(x)=4\cos(x/2)\{x\cos(x/2)-2\sin(x/2)\}.
\]
The expression in braces decreases from zero on $(0,2\pi)$.
Consequently $h$ decreases on $(0,\pi)$ and increases on $(\pi,2\pi)$,
with $h(\pi)=\pi^2-12<0$ and $h(2\pi)=4\pi^2>0$.
Thus $h$ has exactly one zero in $(\pi,2\pi)$. For $x\ge2\pi$,
$r(x)\le1/(4\pi)<3/32\le r(4)$. The zero is therefore the unique global maximizer of $r$.
The scalar checker isolates it in $(3.99,4)$ and encloses
$\kappa_*<0.099161913515<\kappa$ with Arb, a library for rigorous
ball arithmetic.

The cubic expression $V_0t^2/2-\kappa E_0t^3$ is a lower bound for
\eqref{eq:tyurin-exponent} at every frequency. The numerical cosine
lower bound is used only when the entire argument interval lies in
$[4,2\pi]$, which is contained in the cosine branch since $\theta<4$.

\subsection{A logarithmic correction}
For each summand indexed by $J$, define
\[
 h_j=\frac{v_jt^2}{2}-\kappa(b_j+v_j^{3/2})t^3,
 \qquad q_j=\max(h_j,0),\qquad
 H=\frac{V_0t^2}{2}-\kappa(B+\tau_0)t^3.
\]
Since $b_j\ge v_j^{3/2}$,
\begin{equation}\label{eq:q-cap}
 0\le q_j\le m(t,a):=
 \max_{0\le x\le t\sqrt a}(x^2/2-2\kappa x^3)
 \le\frac1{216\kappa^2}<\frac12.
\end{equation}
The maximum in $m$ occurs at $x=\min(t\sqrt a,1/(6\kappa))$.
The single-summand bound gives $|f_{X_j}(t)|\le\sqrt{1-2q_j}$.
Applying $-\tfrac12\log(1-2q)\ge q+q^2$ and using
$\sum_jq_j\ge\sum_jh_j=H$ gives
\begin{equation}\label{eq:log-product}
 \prod_j|f_{X_j}(t)|\le\exp\left\{-H-\sum_jq_j^2\right\}.
\end{equation}

\begin{lemma}\label{lem:quadratic-charge}
Let $A_0=t^2/2-2\kappa\sqrt a\,t^3$. When $A_0\ge0$,
\begin{equation}\label{eq:Q}
 \sum_jq_j^2\ge
 Q:=\left(\max\left\{0,
       A_0\frac{\tau_0}{\sqrt{V_0}}-\kappa t^3(B-\tau_0)
     \right\}\right)^2.
\end{equation}
For $A_0<0$, take $Q=0$.
\end{lemma}
\begin{proof}
Since $v_j\le a$ and $q_j\ge h_j$,
\[
 q_j+\kappa t^3(b_j-v_j^{3/2})
 \ge v_j(t^2/2-2\kappa\sqrt{v_j}\,t^3)\ge A_0v_j.
\]
When $A_0\ge0$, taking $\ell_2$ norms and applying the triangle
inequality gives
\[
 \|(q_j)\|_2\ge A_0\|(v_j)\|_2-\kappa t^3(B-\tau_0).
\]
Cauchy--Schwarz gives $\tau_0\le\sqrt{V_0}\|(v_j)\|_2$.
Taking the maximum of the right side and zero, then squaring, proves
\eqref{eq:Q}.
\end{proof}
After omitting index $i$, the linear sum is $H-h_i\ge H-m$.
Also $\sum_{j\ne i}q_j^2\ge\max(0,Q-m^2)$. Hence
\begin{equation}\label{eq:corrected-products}
 \begin{split}
 F_0(t;V_0,B,a,\tau_0)&=\min\{1,e^{-H-Q}\},\\
 D_0(t;V_0,B,a,\tau_0)&=\min\{1,e^{-H+m-\max(0,Q-m^2)}\}.
 \end{split}
\end{equation}
Here $D_0$ bounds the product after any one index is omitted. We may
also omit the quadratic correction in either expression. When several
upper bounds apply, we take their pointwise minimum. For the empty
index set, set $F_0=D_0=1$.

To bound \eqref{eq:corrected-products} when $\tau_0$ ranges over an
interval $[\ell,h]$, suppose first that $A_0\ge0$ and $t>0$. Put
\[
 k=\kappa t^3,\qquad a_0=A_0/\sqrt{V_0}+k,\qquad c_0=kB.
\]
Apart from terms independent of $\tau_0$, either exponent has the form
\[
 k\tau_0-\max\{0,(a_0\tau_0-c_0)_+^2-m_0^2\},
 \qquad z_+=\max(z,0),
\]
where $m_0=0$ for $F_0$ and $m_0=m$ for $D_0$. Its derivative is $k$
until $a_0\tau_0-c_0=m_0$, and then $k-2a_0(a_0\tau_0-c_0)$.
It is maximized on $[\ell,h]$ at
\[
 \min\{h,\max\{\ell,\tau_*\}\},\qquad
 \tau_*=\frac{c_0+\max\{m_0,k/(2a_0)\}}{a_0}.
\]
When $A_0<0$, $Q=0$ and the exponent is increasing in $\tau_0$.
At $t=0$ both product bounds equal one.

\section{Real and imaginary zero-bias errors}\label{sec:directional}

For a centered variable of positive variance $v$, its zero-bias law is
defined by
\begin{equation}\label{eq:zero-bias}
 v\E g'(X^*)=\E[Xg(X)]
\end{equation}
for absolutely continuous test functions for which the expectations
exist. The first Wasserstein distance between two laws is the infimum
of $\E|U-V|$ over couplings with those laws. By
\citet[Theorem~3]{Tyurin2009}, the distance between $X$ and $X^*$
is at most $\E|X|^3/(2v)$. Applying this bound to $e^{\ii tx}$ gives
\begin{equation}\label{eq:classical-zero-bias}
 |f_X(t)-f_{X^*}(t)|\le |t|\frac{\E|X|^3}{2v}.
\end{equation}
To retain information lost by this bound, let $W$ have mean zero and
variance one, and put $Y=|W|$, $\rho=\E Y^3$, and $e=\rho-1$.
Then $e\ge0$ and
\begin{equation}\label{eq:moment-excess}
 \E[(Y-1)^2(Y+1/2)]=e.
\end{equation}
At $e=0$, $W$ has the Rademacher distribution: it takes the values
$-1$ and $1$ with probability $1/2$ each. Its characteristic function
is $\cos x$, and that of its zero-bias law is $\sinc x$.

\begin{lemma}\label{lem:peano}
For every real $x$,
\begin{equation}\label{eq:real-cf}
 |\Re f_W(x)-\cos x|\le e|x|^3/6,
 \qquad
 |\Re f_{W^*}(x)-\sinc x|\le e|x|/2,
\end{equation}
where $\sinc x=\sin(x)/x$ with its continuous value at zero.
\end{lemma}
\begin{proof}
For an even function $f$ with Lipschitz continuous second derivative,
define
\[
 R_f(y)=f(y)-f(1)-\frac{f'(1)}2(y^2-1),\qquad y\ge0.
\]
Taylor's integral formula gives
$R_f(y)=\int_0^{\max(1,y)}K_y(u)f'''(u)\,du$ with a nonnegative
kernel $K_y$. If $y\le1$, the kernel is $u(y-1)^2/2$ on $[0,y]$ and
$(1-u)(u-y^2)/2$ on $[y,1]$. If $y\ge1$, it is
$u(y-1)^2/2$ on $[0,1]$ and $(y-u)^2/2$ on $[1,y]$.
Its integral is $(y-1)^2(y+1/2)/6$. Apply this identity to
$f(y)=\cos(xy)$, use $\|f'''\|_\infty\le|x|^3$, and take
expectations in \eqref{eq:moment-excess}. The tangent's quadratic term
has expectation zero because $\E Y^2=1$.

Define the square-biased law by
$\E g(W^\square)=\E[W^2g(W)]$ for bounded measurable $g$.
The zero-bias representation is $W^*=UW^\square$, with $U$ uniform
on $[0,1]$ and independent of $W^\square$.
For $z\ge0$,
\[
 \Prob(|W^*|\le z)=\E[Y\min(z,Y)]\le\min(z,1).
\]
Thus $|W^*|$ stochastically dominates $\Unif[0,1]$. Their first
Wasserstein distance is their mean difference
$\rho/2-1/2=e/2$. Apply the Lipschitz bound for $y\mapsto\cos(xy)$.
\end{proof}

\begin{lemma}\label{lem:imaginary}
For all real $x$,
\begin{align}
 |\Im(f_W(x)-f_{W^*}(x))|
   &\le (|x|/2+|x|^3/6)\sqrt{e(\rho+5/3)},\label{eq:imaginary-error}\\
 |\Im f_W(x)|&\le |x|^3\sqrt{e(\rho+5/3)}/6.\label{eq:imaginary-cf}
\end{align}
\end{lemma}
\begin{proof}
It suffices to take $x>0$. Define
$A(y)=x\sinc(xy)+\cos(xy)/x$. The identity \eqref{eq:zero-bias}
and centering give
\[
 \Im(f_W(x)-f_{W^*}(x))=\E[W(A(Y)-A(1))].
\]
Since $|\sinc'(z)|\le |z|/3$, we have
$|A'(y)|\le x(1+x^2/3)y$ and hence
\[
 |A(y)-A(1)|\le (x/2+x^3/6)|y^2-1|.
\]
Cauchy--Schwarz and \eqref{eq:moment-excess} give
\[
 \E[Y|Y^2-1|]\le
 \sqrt{e\,\E\!\left[\frac{Y^2(Y+1)^2}{Y+1/2}\right]}
 \le\sqrt{e(\rho+5/3)}.
\]
The final inequality follows from
\[
 \frac{y^2(y+1)^2}{y+1/2}=y^3+\frac32y^2+
 \frac{y^2}{4(y+1/2)},\qquad
 \frac{y^2}{y+1/2}\le\frac{2+4y^2}{9}.
\]
The second inequality is equivalent to $(y-1)^2(4y+1)\ge0$.
For \eqref{eq:imaginary-cf}, repeat the same argument with
$A(y)=x\sinc(xy)$ and $|A'(y)|\le x^3y/3$.
\end{proof}

\section{A quadratic minorant and the aggregate error}\label{sec:quadratic}

\begin{lemma}\label{lem:supporting-quadratic}
For $0<x\le2$, put
\[
 F_x(y)=\cos(xy)-y\sin(xy)/x,\qquad
 A_x=\cos x+(1+x^2)\sinc x.
\]
Then $F_x(y)\ge F_x(1)-A_x(y^2-1)/2$ for $y\ge0$. Consequently
\begin{equation}\label{eq:one-sided-real}
 \Re(f_W(x)-f_{W^*}(x))\ge\cos x-\sinc x.
\end{equation}
\end{lemma}
\begin{proof}
Set $h_x(z)=\cos z+(1+x^2)\sinc z$ and $A_x=h_x(x)$.
The tangent is $F_x(1)-A_x(y^2-1)/2$. Multiply the difference
from this tangent by $x^2$ and put $z=xy$. Its derivative is
$z[A_x-h_x(z)]$.

The function $h_x$ is strictly decreasing on $(0,\pi)$. On
$[\pi,2\pi]$ it is at most one, and for $z\ge2\pi$ it is at most
$1+5/(2\pi)<11/6$. Alternating Taylor lower bounds give
\[
 A_x\ge p_0(x^2),\qquad
 p_0(u)=2+u/3-7u^2/60+17u^3/2520-u^4/5040.
\]
For $0\le u\le4$,
$p_0''(u)=(-98+17u-u^2)/420<0$, so
$p_0(u)\ge\min(p_0(0),p_0(4))=194/105>11/6$.
The tangent difference decreases up to $z=x$ and increases afterwards.
Its minimum is zero. Taking expectations, and using $\E Y^2=1$,
proves \eqref{eq:one-sided-real}.
\end{proof}

For $0\le x\le2$, let $g=\sinc x-\cos x\ge0$ and
$c=x/2+x^3/6$. Lemmas~\ref{lem:peano} and
\ref{lem:supporting-quadratic} give
\begin{equation}\label{eq:real-interval}
 -g\le\Re(f_W-f_{W^*})\le-g+ce,
 \qquad |\Re(f_W-f_{W^*})|^2\le g^2+c^2e^2.
\end{equation}
The alternating series gives the increasing upper bound
\begin{equation}\label{eq:alpha}
 \frac{2g}{x}\le\alpha(x):=
 \begin{cases}
  2x/3-x^3/15+x^5/420,&0\le x\le2,\\
  2x/3,&x>2.
 \end{cases}
\end{equation}
The quotient at zero is defined by continuity. On $[0,2]$ the
polynomial is at most $2x/3$. Put $q(x)=1+x^2/3$.
Combining \eqref{eq:real-interval} with \eqref{eq:imaginary-error} and
\eqref{eq:classical-zero-bias}, we obtain
\begin{equation}\label{eq:directional-prefactor}
 |f_W(x)-f_{W^*}(x)|\le\frac{|x|}{2}
 \min\left\{\rho,
   \sqrt{\alpha(|x|)^2+q(x)^2(\rho-1)(2\rho+2/3)}\right\}.
\end{equation}
For $x>2$ the square root is at least $\rho$, so the minimum simply
uses \eqref{eq:classical-zero-bias}. Reflection and continuity handle
negative $x$ and zero.

\begin{proposition}\label{prop:aggregation}
For a nonempty subset of summands with moment sums $B,\tau_0$ and
individual variances at most $a$,
\begin{equation}\label{eq:aggregate-error}
 \sum_j v_j|f_{X_j}(t)-f_{X_j^*}(t)|\le\frac t2
 P(t;B,a,\tau_0),\qquad t\ge0,
\end{equation}
where, with $x=t\sqrt a$,
\begin{equation}\label{eq:P}
 P(t;B,a,\tau_0)=\min\left\{B,
  \sqrt{\alpha(x)^2\tau_0^2+q(x)^2(B-\tau_0)(2B+2\tau_0/3)}\right\}.
\end{equation}
The bound $P$ increases with $t,a,B$ and decreases with $\tau_0$.
For the empty subset, set $P=0$.
\end{proposition}
\begin{proof}
The square root in \eqref{eq:directional-prefactor} is increasing in
$x$ and concave in $\rho\ge1$. In the variable $e=\rho-1$, its
quadratic has discriminant $(64/9)q^4-8q^2\alpha^2>0$, since
$\alpha\le2x/3$ and
$q^2-4x^2/3=(1-x^2/3)^2\ge0$.
The second derivative of the square root has the opposite sign to this
discriminant; continuity handles a zero radicand at an endpoint.
After scaling \eqref{eq:directional-prefactor} to variance $v_j$,
each term has the factor $tv_j^{3/2}/2$ and moment ratio
$\rho_j=b_j/v_j^{3/2}$. Replace $t\sqrt{v_j}$ by $t\sqrt a$ and
apply Jensen's inequality with weights $v_j^{3/2}/\tau_0$; their
weighted mean of $\rho_j$ is $B/\tau_0$. This gives the square root
in \eqref{eq:P}. The classical bounds also give $tB/2$.

The expression under the square root in \eqref{eq:P} is
\[
 2q^2B^2-\frac43q^2B\tau_0+
 (\alpha^2-2q^2/3)\tau_0^2.
\]
Its nonconstant coefficients in $\tau_0$ are nonpositive. Its other
monotonicities follow from the original nonnegative-factor expression.
\end{proof}

\subsection{A characteristic-function bound using moment excess}
Squaring the real-part estimate in Lemma~\ref{lem:peano} and adding
the square of the bound in \eqref{eq:imaginary-cf} gives another
single-summand bound. If
$\tau_j=v^{3/2}$ and $C=|\cos(t\sqrt v)|$, then
\begin{equation}\label{eq:paired-cf}
 |f_X(t)|^2\le C^2+\frac{C(b-\tau_j)|t|^3}{3}
       +\frac{t^6(b-\tau_j)(2b+2\tau_j/3)}{36}.
\end{equation}
This holds for every real $t$. We use the smaller of its square root
and \eqref{eq:single-cf}, evaluated at $|t|$.
The displayed terms increase with $C,b,|t|$ and decrease with $\tau_j$
when $b\ge\tau_j$. On a parameter interval, first bound $C$ over the
entire cosine argument interval, then apply these monotonicities.

\section{A summand of maximal variance}\label{sec:distinguished}

The independent-sum zero-bias identity, with total variance one, is
\begin{equation}\label{eq:sum-zero-bias}
 f_{S^*}(t)=\sum_i v_i f_{X_i^*}(t)\prod_{j\ne i}f_{X_j}(t).
\end{equation}
For the remaining summands, use the product bound
\[
 R(t)=\min\{e^{-\psi(t;w,B_R+\tau_R)},
                  F_0(t;w,B_R,a_R,\tau_R)\}.
\]
If the remainder is empty, set $R(t)=D_R(t)=1$.
Let $A(t)$ be the minimum of the single-summand bounds
\eqref{eq:single-cf} and the square root of \eqref{eq:paired-cf} for
$X_1$. Let $D_R(t)$ be the bound $D_0$ for the remaining summands,
and $D_S(t)$ the corresponding bound for the whole sum. We also evaluate
these bounds with $Q$ omitted and take the smaller interval upper bounds.
Proposition~\ref{prop:aggregation} gives

\begin{equation}\label{eq:distinguished-error}
 \begin{split}
 |f_S(t)-f_{S^*}(t)|\le\frac t2\bigl\{&
 P(t;b,d,d^{3/2})R(t)\\
 &+P(t;B_R,a_R,\tau_R)A(t)D_R(t)\bigr\}.
 \end{split}
\end{equation}
The whole-sum estimate is $(t/2)P(t;L,d,\tau)D_S(t)$, and the
classical estimate is $tL/2$. Thus
$|f_S(t)-f_{S^*}(t)|\le(tL/2)\mathcal B(t)$, where
\begin{equation}\label{eq:B}
 \mathcal B(t)=\min\left\{1,\frac{P(t;L,d,\tau)D_S(t)}L,
 \frac{P(t;b,d,d^{3/2})R(t)+P(t;B_R,a_R,\tau_R)A(t)D_R(t)}L
 \right\}.
\end{equation}
This formula also defines the endpoint value $\mathcal B(0)$.

For the characteristic function itself, use
\begin{equation}\label{eq:whole-cf}
 F_S(t)=\min\{e^{-\psi(t;1,L+\tau)},\ A(t)R(t),\
                    F_0(t;1,L,d,\tau)\}.
\end{equation}

All terms use the same actual $\tau$; an interval evaluation allows its
whole range in the parameter box. The calculation also bounds the last
term of \eqref{eq:B} by separating the moment fraction
$\lambda=b/L$ from two normalized errors. Put
\[
 U_1=\frac{P(t;b,d,d^{3/2})R(t)}b,\qquad
 U_R=\frac{P(t;B_R,a_R,\tau_R)A(t)D_R(t)}{B_R}.
\]
If $B_R=0$, set $U_R=0$. Both quantities lie in $[0,1]$, and the
last term in \eqref{eq:B} equals $\lambda U_1+(1-\lambda)U_R$.
Suppose a box gives $\lambda\in[\lambda_-,\lambda_+]\subseteq[0,1]$
and $U_i\le\overline U_i$. Then
\begin{equation}\label{eq:moment-fraction}
 \lambda U_1+(1-\lambda)U_R\le
 \max_{\lambda\in\{\lambda_-,\lambda_+\}}
       \{\lambda\overline U_1+(1-\lambda)\overline U_R\}.
\end{equation}
The right side is the maximum of an affine function on an interval.
We take the smaller of this bound and a direct interval evaluation of
the numerator in \eqref{eq:B}, divided by a lower bound for $L$.
When a denominator in the separate bounds for $U_i$ has lower endpoint
zero, the bound $U_i\le1$ is used.

Write $D(t)=f_S(t)-e^{-t^2/2}$. The zero-bias identity gives
$f_S'(t)=-t f_{S^*}(t)$, and hence
\[
 D'(t)+tD(t)=t(f_S(t)-f_{S^*}(t)),\qquad D(0)=0.
\]
Multiplication by $e^{t^2/2}$ and integration yield
\[
 D(t)=e^{-t^2/2}\int_0^t u(f_S(u)-f_{S^*}(u))e^{u^2/2}\,du.
\]
Taking absolute values gives
\begin{equation}\label{eq:ode-comparison}
 \frac{|f_S(t)-e^{-t^2/2}|}{L}
 \le R_\Delta(t):=e^{-t^2/2}
       \int_0^t \frac{u^2}{2}\mathcal B(u)e^{u^2/2}\,du.
\end{equation}

\section{Smoothing and numerical integration}\label{sec:smoothing}

For $0<u<1$, put
\[
 K(u)=\frac{1-u}{2}+
 \frac{\ii}{2}\left((1-u)\cot(\pi u)+\frac1\pi\right).
\]
Prawitz's inequality, in the form stated by
\citet[Lemma~4.1]{Shevtsova2012} and \citet[eq.~(46)]{Tyurin2009},
gives for any $T>0$ and $0<s<1$
\begin{equation}\label{eq:prawitz}
 \frac\Delta L\le
 2\int_0^s|K(u)|R_\Delta(Tu)\,du+
 \frac{2\int_s^1|K(u)|F_S(Tu)\,du+G(T,s)}{L},
\end{equation}
where
\begin{equation}\label{eq:gaussian-correction}
 G(T,s)=2\int_0^s\left|K(u)-\frac{\ii}{2\pi u}\right|
          e^{-T^2u^2/2}\,du+\frac{E_1(T^2s^2/2)}{2\pi},
 \qquad E_1(z)=\int_z^\infty\frac{e^{-v}}v\,dv.
\end{equation}

\subsection{Integration over a frequency interval}
On a frequency interval $[a,b]$, let $\mathcal B_j$ be an upper bound for
$\mathcal B$, and write
\[
 I_a=\int_0^a u^2\mathcal B(u)e^{u^2/2}/2\,du.
\]
Let $A_j$ be an upper bound for $2\pi u|K(u)|$ on
$u\in[a/T,b/T]$. Under the change of variable $t=Tu$, the first term of
\eqref{eq:prawitz}, restricted to $[a,b]$, is at most
\begin{equation}\label{eq:fubini-cell}
 \frac{A_j}{\pi}\left\{
 I_a\frac{E_1(a^2/2)-E_1(b^2/2)}2+
 \mathcal B_j\left(\frac{b^3-a^3}{18}
                  -\frac{a^3}{6}\log\frac ba\right)\right\}.
\end{equation}
To see this, split the integral defining $R_\Delta(t)$ at $a$.
The part up to $a$ contributes
$\int_a^b I_a e^{-t^2/2}/t\,dt$, which equals the first term in
braces. For the remaining part, use $e^{-(t^2-u^2)/2}\le1$ and
interchange the nonnegative integrals:
\[
 \int_a^b\frac1t\int_a^t\frac{u^2}{2}\,du\,dt
 =\frac{b^3-a^3}{18}-\frac{a^3}{6}\log\frac ba.
\]
At $a=0$, $I_a=0$ and the expression in braces in
\eqref{eq:fubini-cell} is $\mathcal B_jb^3/18$.
In the calculation, $I_a$ is replaced by an upper bound accumulated
from the preceding intervals. For $a>0$, we also integrate the bound
$|f_S-e^{-t^2/2}|\le F_S+e^{-t^2/2}$ over $[a,b]$ and use the
smaller of the two upper bounds for that integral.

\subsection{The Gaussian correction}
Direct subtraction of the singular part of $K$ gives
\begin{equation}\label{eq:regular-kernel}
 2\left|K(u)-\frac{\ii}{2\pi u}\right|
   =(1-u)\sqrt{1+h(u)^2},\qquad
 h(u)=\frac1{\pi u}-\cot(\pi u).
\end{equation}
Set $h(0)=0$ by continuity. The partial-fraction identity \citep[eq.~4.22.3]{DLMF} is
\[
 h(u)=\frac{2u}{\pi}\sum_{n=1}^\infty\frac1{n^2-u^2}.
\]
For $0\le u\le s<1$ and any integer $M\ge1$, define
\[
 c_M(s)=\frac{\pi^2/6-\sum_{n=1}^M n^{-2}}
                  {1-s^2/(M+1)^2},\qquad
 h_M(u)=\frac{2u}{\pi}\left(\sum_{n=1}^M\frac1{n^2-u^2}
                                   +c_M(s)\right).
\]
For $n\ge M+1$ and $0\le u\le s$,
$1/(n^2-u^2)\le n^{-2}/(1-s^2/(M+1)^2)$. Summing gives
$0\le h(u)\le h_M(u)$. In the computation, $c_M(s)$ is replaced by
an upper dyadic endpoint. Since $s<1$, every denominator in the finite
sum is positive on $[0,s]$. Hence
\begin{equation}\label{eq:gaussian-majorant}
 G(T,s)\le \int_0^s(1-u)\sqrt{1+h_M(u)^2}e^{-T^2u^2/2}\,du
       +\frac{E_1(T^2s^2/2)}{2\pi}.
\end{equation}
The calculation uses $M=8$ and Arb to enclose the integral in
\eqref{eq:gaussian-majorant}; the integration routine is described in
Section~\ref{sec:arithmetic}. We also evaluate the following bound and
use the smaller upper endpoint:

\begin{equation}\label{eq:simple-gaussian}
 G(T,s)\le\int_0^s(1-u+\pi^2u^2/18)e^{-T^2u^2/2}\,du
               +\frac{E_1(T^2s^2/2)}{2\pi},
\end{equation}
also supplied by \citet[Lemma~4.1]{Shevtsova2012}. Its integral is
expressible using the error function and elementary functions.

\section{The interval calculation}\label{sec:certificate}

The interval calculation covers $L\in[L_-,L_+]$, where
$L_-\approx0.006$ and $L_+\approx1.15$ are the following exact dyadic
rationals, represented by the stored binary64 endpoints:

\begin{equation}\label{eq:exact-endpoints}
 L_-=\frac{3458764513820541}{576460752303423488},\qquad
 L_+=\frac{2589569785738035}{2251799813685248}.
\end{equation}

\subsection{Small and large moments}
For $0<L\le L_-$, Corollary~4.18 of \citet{Shevtsova2012} gives
\begin{equation}\label{eq:small-L}
 \Delta/L\le\CE+0.3413L^{1/3}
      \le\CE+0.3413L_-^{1/3}<\SmallUpper.
\end{equation}
This corollary applies because $L_-<0.01$. The scalar checker
recomputes the last enclosure.

For every mean-zero variance-one law, Cantelli's inequality and
reflection give
\begin{equation}\label{eq:cantelli}
 \Delta\le\sup_{x>0}\left(\Phi(x)-\frac{x^2}{1+x^2}\right)
      <\VarianceUpper.
\end{equation}
For completeness, at $x\ge0$ the lower CDF bound follows from Cantelli,
while the opposite difference is at most $1-\Phi(x)\le1/2$.
Reflection handles $x<0$, including either one-sided convention at atoms.
The displayed supremum exceeds $1/2$. Its stationary point is unique:
the logarithmic derivative of
$2x/((1+x^2)^2\phi(x))$ is
$(x^2-1)^2/(x(1+x^2))\ge0$, and the ratio increases from zero to
infinity. Arb isolation of the stationary point gives
\begin{equation}\label{eq:large-L}
 \Delta/L<\LargeUpper\qquad(L\ge L_+).
\end{equation}

\subsection{Subdivision of the parameter domain}
For each closed interval $L\in[\ell,H]$, the calculation starts from
\[
 0\le d\le\min\{1,H^{2/3}\},\qquad
 0\le b\le H,\qquad 0\le\tau\le\min\{1,H\},
\]
with outward-rounded endpoints. The constraints in
\eqref{eq:moment-domain} tighten this box. A box is accepted when its
computed upper bound is finite and below $0.475$ throughout the box.
It is discarded when the constraints prove it infeasible; otherwise,
it is divided into two closed boxes whose union contains it. The stored
tree records each split and terminal decision. The finer claimed
constant $\CandidateUpper$ is checked against the recorded maxima
after this subdivision is complete.

For each box, the evaluator tries three pairs of smoothing parameters
$T,s$ and takes the minimum of the resulting bounds. Each pair divides
$[0,Ts]$ and $[Ts,T]$ into 512 equal frequency intervals.
The pairs are chosen by the rules in
\path{verify/kernel/refined_integral.py}, with scale and shift pairs
$(1,0)$, $(1.02,-0.02)$, and $(1.04,-0.02)$ from
\path{verify/kernel/batch_weights.py}. The upper endpoint of the
box's $\tau$ interval selects an input to these rules through
\path{verify/kernel/gaussian_batch.py}. The resulting binary64 values
of $T,s$ are treated as exact dyadic parameters; each bound is evaluated
over the actual moment box.

Before subdivision, the calculation tries two bounds uniform over
$(d,b,\tau)$. The first uses \eqref{eq:cantelli} in the looser form
$\Delta/L\le0.54093655/\ell$. The second applies the smoothing
integral with $d_U=H^{2/3}$ and $E_U=2H$, rounded upward, using
\[
 \mathcal B_U(t)=\min\{1,e^{-t^2/2+\kappa E_Ut^3+m(t,d_U)}\},
 \qquad F_U(t)=e^{-\psi(t;1,E_U)}.
\]
These follow from $\tau\le L\le H$ and the product bounds of
Section~\ref{sec:products}. If either proves the acceptance bound for
the whole $L$ interval, the record has a single accepted node.

\begin{center}
\begin{tabular}{lr}
\toprule
Quantity & Value\\
\midrule
Closed $L$ intervals & \BandCount\\
Partition nodes & \NodeCount\\
Accepted leaves & \AcceptedCount\\
Infeasible leaves & \InfeasibleCount\\
Largest recorded upper endpoint & \RecordedUpper\\
Upper bound for $L\le L_-$ & $<\SmallUpper$\\
Upper bound for $L\ge L_+$ & $<\LargeUpper$\\
\bottomrule
\end{tabular}
\end{center}
Of the \BandCount{} intervals, 545 used the moment-parameter bounds, 215 used the
uniform smoothing bound, and seven used the variance bound. Every
accepted leaf has a finite upper endpoint. Together with the two
complementary estimates, their maximum gives
\[
 \Cind\le\max\{\RecordedUpper,\SmallUpper,\LargeUpper\}
 <\CandidateUpper.
\]

\subsection{Completed computation and subsequent checks}
The original calculation evaluated every accepted-leaf bound. The
supplement records three kinds of subsequent checks:
\begin{enumerate}[leftmargin=*,itemsep=3pt]
\item \emph{Saved-record checking} checks file and source-code digests,
tree syntax and counts, contiguous $L$ intervals, recorded upper
endpoints, and the scalar enclosures in
\eqref{eq:small-L}--\eqref{eq:large-L}. It does not recompute the
quadrature at accepted leaves.
\item \emph{Geometric replay} uses the archived implementation to
reconstruct every box reduction and split, check every infeasibility
decision, and verify coverage of all \BandCount{} intervals. It also
checks the analytic complements.
\item \emph{Numerical replay} additionally recomputes quadrature at
accepted leaves. The recorded replay recomputed all 4,944 accepted
leaves in three intervals, selected by lower endpoints nearest $0.006$,
$0.25$, and $0.48$. The checksums of their ordered numerical results and
their maxima matched the original records. Quadrature in the remaining
intervals has not all been recomputed.
\end{enumerate}
The records contain these checksums and maxima; individual numerical
values can be recovered by repeating the computation. Independent
mathematical review remains outstanding.

\subsection{Arithmetic assumptions and reproduction}\label{sec:arithmetic}
The transcendental constants are enclosed with 100-bit Arb arithmetic.
The vector calculation assumes IEEE-754 binary64 elementary operations
and square root, rounding to nearest, gradual underflow, and no unsafe
fast-math transformations. Interval operations use \texttt{nextafter}
to round endpoints outward. Nonnegative cumulative sums and dot products
are multiplied by $1+4N2^{-53}$ and rounded upward, where $N=512$ is
the number of frequency intervals in each range for each smoothing
choice. The full factor is used for every prefix of a cumulative sum.

For a nonnegative lower bound $p$ on an exponent, the negative
exponential is bounded using
\[
 P_{16}(x)=\sum_{j=0}^{16}\frac{x^j}{j!},\qquad
 e^{-p}\le P_{16}(p/64)^{-64}\quad(0\le p<25).
\]
The implementation evaluates this polynomial with downward rounding,
takes the reciprocal upward, and performs six upward-rounded squarings.
For $p\ge25$, it uses an upper enclosure of $e^{-25}$.
The routine is \texttt{exp\_upper\_from\_nonnegative\_lower} in
\path{verify/kernel/interval_bounds.py}.

The Gaussian integral in \eqref{eq:gaussian-majorant} is evaluated by
\texttt{gaussian\_G} in \path{verify/kernel/gaussian_batch.py}, using
Arb's \texttt{acb.integral}. Its callback passes the integration
routine's analyticity requirement to the square root through
\texttt{sqrt(analytic=analytic)}. The returned ball must be finite and
its imaginary part must contain zero. The validity of the integration
also depends on Arb's integration contract and the callback's branch
handling. The summation allowance, exponential bound, and integration
procedure are included in the outstanding arithmetic audit.

The nine arithmetic source files are identified by SHA-256 digests.
The recorded numerical runtime is Python~3.12.10, NumPy~2.4.6,
python-flint~0.8.0, and SciPy~1.15.3. The complete interval records are
in \path{certificates/cover.json}; the proposed bound and scalar
endpoints are summarized in \path{result.json}.
The supplement's \path{verify/README.md} gives the commands for checking
the saved records, reconstructing the partition, and recomputing
selected numerical bounds.

\section{Possible refinements}\label{sec:questions}

Two relaxations remain in the estimates: the triangle inequality
replaces the modulus of a sum of complex zero-bias errors by the sum of
their moduli, and the moment constraints may admit unrealizable tuples.
A finer parameter partition can reduce interval-enclosure error, but
cannot recover the information lost in either relaxation.

\subsection{Higher-order terms in the logarithm}
The following higher-order bounds were not used in the completed
calculation. With the notation of Lemma~\ref{lem:quadratic-charge},
H\"older interpolation gives, for every integer $p\ge2$,
\[
 \|(v_j)\|_{3/2}
 \le \|(v_j)\|_1^{(2p-3)/(3p-3)}\|(v_j)\|_p^{p/(3p-3)}.
\]
Since $\|(v_j)\|_{3/2}=\tau_0^{2/3}$ and $\|(v_j)\|_1=V_0$,
\[
 \|(v_j)\|_p\ge\frac{\tau_0^{2-2/p}}{V_0^{2-3/p}}.
\]

If $A_0\ge0$, the same norm triangle inequality gives
\begin{equation}\label{eq:higher-charge}
 \sum_jq_j^p\ge Q_p:=
 \left(\max\left\{0,A_0\frac{\tau_0^{2-2/p}}{V_0^{2-3/p}}
                      -\kappa t^3(B-\tau_0)\right\}\right)^p.
\end{equation}
Take $Q_p=0$ otherwise. For $0\le q<1/2$,
\[
 -\tfrac12\log(1-2q)=\sum_{p=1}^{\infty}\frac{2^{p-1}}p q^p.
\]
Truncating the series and using \eqref{eq:higher-charge} gives the
full-product bound below. Omitting an index reduces each $p$th-power
sum by at most $m^p$, which gives the second bound. For any integer
$P\ge2$,
\begin{equation}\label{eq:higher-products}
 \begin{split}
 \prod_j|f_{X_j}(t)|
   &\le\exp\left\{-H-\sum_{p=2}^P\frac{2^{p-1}}p Q_p\right\},\\
 \prod_{j\ne i}|f_{X_j}(t)|
   &\le\exp\left\{-H+m-\sum_{p=2}^P\frac{2^{p-1}}p
                              \max(0,Q_p-m^p)\right\}.
 \end{split}
\end{equation}
Take the minimum of each displayed bound, one, and any other applicable
product bound.

\subsection{Aggregation and phase information}
The real-part interval in \eqref{eq:real-interval} gives
\[
 |\Re(f_W-f_{W^*})|^2\le\max\{g^2,(ce-g)^2\}.
\]
This improves the single-summand bound $g^2+c^2e^2$. Aggregating the
sharper expression by Jensen's inequality requires a concavity argument
for the resulting function of the moments. Such an argument is not
supplied here. Another possibility is to retain relative Fourier phases
for several summands while bounding the remaining errors uniformly.

\clearpage
\bibliographystyle{plainnat}
\bibliography{refs}
\end{document}
